\begin{english}
\chapter*{Abstract}
\addcontentsline{toc}{chapter}{Abstract}
Non-fungible tokens (NFTs) established the notion of a unique digital asset on public
blockchains, yet their markets remain organised almost exclusively around sales for
cryptocurrency: fixed prices, auctions and offers. Direct NFT-for-NFT exchange, especially
between items of different collections, is not served by existing platforms and, where it is
attempted at all, requires trusting the counterparty or an intermediary.

This diploma thesis designs and implements a barter marketplace on the Sui blockchain, in
which a user submits a bid consisting of one or more NFTs from any collection and,
optionally, an amount of SUI, in exchange for a specific NFT held by another user. The
offered assets are escrowed inside a shared \texttt{Bid} object at submission time, so that
acceptance by the owner of the requested NFT completes the swap atomically and without a
custodian. The smart contract is written in Move and exploits Sui's object-centric model.

Barter was historically the first form of economic exchange, long before money was
established as a common measure of value, and it never ceased to have a place where goods
are unique and their value subjective. The digitisation of real-world goods as NFTs has
already moved from the experimental to the institutional stage, a notable example being the
tokenisation of real-estate title deeds on a blockchain by the Dubai Land Department. In this context, a trustworthy
mechanism for direct exchange, one that does not impose the liquidation of every asset into
cryptocurrency as an intermediate step, ceases to be a complement to existing markets and
becomes necessary infrastructure. This thesis proposes such a mechanism, in which the swap is
executed atomically by the contract itself, so that neither party needs to trust the other.

\vspace{1cm}
\noindent\textbf{Keywords:} blockchain, Sui, Move, NFT, marketplace, barter, bid, escrow,
atomic swap, smart contract, dApp, cross-collection
\end{english}
\clearpage
